\documentclass[UTF8,twocolumn,10pt,a4paper]{ctexart}

\usepackage[a4paper,left=1.55cm,right=1.55cm,top=1.45cm,bottom=1.60cm,columnsep=0.66cm]{geometry}
\usepackage{fontspec}
\setmainfont{DejaVu Serif}
\setsansfont{DejaVu Sans}
\setmonofont{DejaVu Sans Mono}
\IfFontExistsTF{Noto Serif CJK SC}{\setCJKmainfont{Noto Serif CJK SC}}{}
\IfFontExistsTF{Noto Sans CJK SC}{\setCJKsansfont{Noto Sans CJK SC}}{}
\IfFontExistsTF{Noto Sans Mono CJK SC}{\setCJKmonofont{Noto Sans Mono CJK SC}}{}

\usepackage{amsmath,amssymb,mathtools,bm}
\usepackage{booktabs,tabularx,array,multirow}
\usepackage{graphicx}
\usepackage{enumitem}
\usepackage{microtype}
\usepackage{xcolor}
\usepackage{balance}
\usepackage{fvextra}
\usepackage[hidelinks]{hyperref}

\hypersetup{
  pdftitle={TGVF Policy RL Current Run Report: Crop Reproduction and TGVF Reward Plan},
  pdfauthor={TGVF Project},
  pdfsubject={Current Qwen3-VL-Instruct policy-RL configuration, T1 data preparation, Crop and TGVF protocols, rewards, judges, and evaluation plan},
  pdfkeywords={TGVF, Crop, GRPO, Qwen3-VL-Instruct, DeepEyes, data selection, reward, judge}
}

\setlength{\parindent}{1em}
\setlength{\parskip}{0.12em}
\renewcommand{\baselinestretch}{1.015}
\setlength{\abovedisplayskip}{4.4pt plus 1pt minus 1pt}
\setlength{\belowdisplayskip}{4.4pt plus 1pt minus 1pt}
\setlength{\abovedisplayshortskip}{2.4pt plus 1pt}
\setlength{\belowdisplayshortskip}{2.4pt plus 1pt minus 1pt}
\setlength{\textfloatsep}{7pt plus 2pt minus 2pt}
\setlength{\dbltextfloatsep}{7pt plus 2pt minus 2pt}
\setlength{\intextsep}{6pt plus 2pt minus 2pt}
\setlength{\emergencystretch}{1.5em}
\ctexset{
  section={beforeskip=0.72ex plus .2ex,afterskip=0.34ex plus .1ex,format=\large\bfseries},
  subsection={beforeskip=0.58ex plus .2ex,afterskip=0.26ex plus .1ex,format=\normalsize\bfseries},
  paragraph={beforeskip=0.45ex plus .15ex,afterskip=0.55em,runin=true,format=\normalsize\bfseries}
}
\setlist{nosep,leftmargin=1.45em}
\allowdisplaybreaks

\newcolumntype{Y}{>{\raggedright\arraybackslash}X}
\newcolumntype{C}{>{\centering\arraybackslash}X}
\newcommand{\TGVF}{\textsf{TGVF}}
\newcommand{\code}[1]{\texttt{#1}}
\newcommand{\pp}{\,\mathrm{pp}}
\newcommand{\Rcrop}{R_{\mathrm{crop}}}
\newcommand{\Rtgvf}{R_{\mathrm{tgvf}}}

\title{\vspace{-1.0em}\textbf{TGVF Policy RL 当前实验报告\newline
Instruct 数据筛选、Crop 复现与 TGVF Reward 路线}}
\author{TGVF Project}
\date{运行快照、实验合同与后续计划，2026 年 8 月 8 日}

\begin{document}
\maketitle
\vspace{-1.55em}

\begin{abstract}
本文冻结当前 Policy RL 阶段的真实实验状态，并把 Crop 与 \TGVF{} 两条线分开记录。当前正在运行的不是 RP66/\TGVF{}，而是基于 Qwen3-VL-8B-Instruct 的 PRL-13-A native Crop/DeepEyes fidelity efficacy gate：8 张 B200、全参数训练（vision/projector/language 均可训练）、GRPO、prompt batch 16、每 prompt 16 条 rollout，即每个 optimizer step 256 条 trajectory；本次只跑到 step 8，不是 run ID 中名义的 BS256$\times$80 正式训练。

训练数据来自 Instruct 模型的完整 T1 筛选：271,842 个候选各用原图无工具生成 8 次，保留天然正确次数在 1--7 之间的样本，最终得到 77,541 条 VStar/ArxivQA/ThinkLite 混合数据。当前 Crop visual reward 为 $0.8A+0.2F+1.2C$；每条 VStar/ArxivQA trajectory 的 $A$ 均由 text-only Qwen2.5-72B-Instruct 判定，Crop 工具自身不使用 VLM judge。ThinkLite 使用 $1.2A+0.4F$，先走规则 math verifier，必要时使用同一 72B fallback。

报告同时给出未来 \TGVF{} arm 的 prompt、RP66 representation 绑定和 Stage3-shaped reward：$2A_{\rm gated}+T+F_{\rm focus}+G_{\rm grounding}+P$。其中答案判定仍是 rule-first/Qwen2.5-72B；只有 $F/G$ 需要 Qwen3-VL-32B-Thinking 视觉 judge。已完成的 PRL-09-R3 因关闭 $F/G$ 而没有加载 32B，因此不能写成“完整 shaped reward 已验证有效”。截至本文快照，当前 Crop 小门控运行数值稳定、无 OOM/aborted/judge failure，但 on-policy batch reward 不能替代固定 step-0 对 step-8 benchmark；是否提升仍须等待配对评测。
\end{abstract}

\noindent\textbf{关键词：} Policy RL；Crop；TGVF；GRPO；Qwen3-VL-Instruct；T1 数据筛选；reward judge

\section{首页结论与范围}

\begin{figure*}[t]
\centering
\fbox{\parbox{0.965\textwidth}{\small
\textbf{Current experiment contract.}
当前活动 run 是 \textbf{PRL-13-A Crop efficacy gate}，不是 \TGVF{} run。基础模型为 Qwen3-VL-8B-Instruct；当前 override 为 BS16$\times$n16$\times$8 steps；8 GPU；full tuning；LR $10^{-6}$。名义 parent contract 才是 BS256$\times$n16$\times$80，不能把两者混写。VStar/ArxivQA 的 final-answer correctness 使用 \textbf{text-only Qwen2.5-72B-Instruct}，没有加载 32B VLM。\TGVF{} 里的 \textbf{Qwen3-VL-32B-Thinking} 只属于完整 focus/grounding quality reward，当前 Crop 不使用它。当前训练 reward、工具率和梯度均为健康性证据；只有固定 step-0/step-8 evaluation 才能回答“是否提升准确率”。}}
\caption{本文最重要的身份边界。}
\label{fig:contract}
\end{figure*}

本文回答四个问题：当前到底在跑什么；训练数据如何得到；Crop 和 \TGVF{} 分别看到什么 prompt、得到什么 reward；本次小门控通过后如何做可扩展验证。历史 Thinking-model、decoder-LoRA、错误 Crop 坐标和 RP66 小 batch 结果只作为诊断背景，不与当前 full-tuning Crop 直接混合比较。

\begin{table}[t]
\centering
\caption{两条方法线与当前状态。}
\label{tab:lines}
\small
\begin{tabularx}{\columnwidth}{lYY}
\toprule
线 & 当前角色 & 状态 \\
\midrule
Crop & DeepEyes positive control；验证 full vision-policy RL 是否能在小规模产生正确方向 & PRL-13-A step-8 gate 正在/刚完成运行，待固定评测 \\
\TGVF{} & 学习何时请求 target-conditioned $D$，并利用 image+$D$ & 旧小 LoRA 与 PRL-09 已完成；full-scale matched arm 尚未启动 \\
Representation & 为 \TGVF{} 工具提供冻结 Adapter & RP66 已证明 additive utility；RP67 只宜作独立表示消融 \\
\bottomrule
\end{tabularx}
\end{table}

\section{当前 PRL-13-A 的真实配置}

\subsection{运行身份}

当前 launcher 为 \path{/tmp/launch_prl13_formal_micro32_to8.py}；输出目录简称 \code{PRL13/pilot-bs16-n16-micro32}。W\&B run ID 为 \code{g9j7mkru}，可点击查看
\href{https://wandb.ai/mio_nora/tgvf-policy-rl/runs/g9j7mkru}{online run g9j7mkru}。

长 run ID 含 \code{BS256-N16-...-80STEP}，但 launcher 在 compose 后覆盖了 batch、world size、micro batch 和 target step。表~\ref{tab:parent-live} 是必须用于解释结果的最终口径。

\begin{table}[t]
\centering
\caption{名义 parent contract 与本次 live override。}
\label{tab:parent-live}
\scriptsize
\begin{tabularx}{\columnwidth}{Ycc}
\toprule
项目 & Parent formal & 当前 gate \\
\midrule
prompt batch/step & 256 & 16 \\
rollout/prompt & 16 & 16 \\
trajectory/step & 4,096 & 256 \\
optimizer steps & 80 & 8 \\
总 prompts & 20,480 & 128 \\
总 trajectories & 327,680 & 2,048 \\
GPU/world size & 4 & 8 \\
actor micro/GPU & 4 & 32 \\
FSDP/offload & 4-way, offload & 8-way, no offload \\
save & 1/8/20/45/80 & 每 step，保留最新 actor \\
\bottomrule
\end{tabularx}
\end{table}

\subsection{模型、优化器与 rollout}

\begin{table}[t]
\centering
\caption{当前 gate 的 live resolved hyperparameters。}
\label{tab:live-hp}
\scriptsize
\begin{tabularx}{\columnwidth}{lY}
\toprule
项目 & 实际值 \\
\midrule
Policy & Qwen3-VL-8B-Instruct \\
训练范围 & full；LoRA rank 0；vision/projector/language 全部更新 \\
算法 & GRPO，$n=16$，组内 advantage 按标准差归一化 \\
Optimizer & AdamW；LR $10^{-6}$ constant；$\beta=(0.9,0.999)$；WD 0.01 \\
Update & PPO epoch 1；clip $\epsilon=0.2$；dual clip 3；grad clip 1 \\
Regularizer & KL reward 0；actor KL loss off；entropy coefficient 0 \\
长度 & prompt 8,192；response 20,480；model len 32,768 \\
Sampling & vLLM async；$T=1$，top-$p=1$，top-$k=-1$；TP=1 \\
Agent & 8 workers；最多 7 assistant / 6 observation turns；单 turn 一次工具调用 \\
视觉 & 原图加至多 6 张 crop；Qwen $0..1000$ bbox 映射至 immutable source pixels \\
并行 & 8$\times$B200，FSDP2；BF16 compute；fused kernels；text FlexAttention、vision SDPA \\
Checkpoint & save every step；step 8 内置 T1-PROBE256；无自动 CoreDev \\
\bottomrule
\end{tabularx}
\end{table}

这次 full tuning 是相对旧 pilot 的最大结构差异之一：Crop observation 重新经过原生 vision encoder，而 vision tower 本身也可被 reward 更新。旧 decoder-LoRA 设置只能学习“如何调用/读取已有视觉表示”，不能改变视觉编码路径，因此其失败不能直接推出 full Crop 也失败。

\section{Instruct T1 数据制备}

\subsection{候选池与生成合同}

T1 selector 使用同一个 Qwen3-VL-8B-Instruct，而不是 Thinking model。三个已做 CoreDev exact image-byte leakage screening 的来源文件合计 361,709 rows；再以 outcome-independent content-hash bottom-$k$ 取固定 source quota：VStar 170,000、ArxivQA 32,000、ThinkLite 69,842，共 271,842 candidates。

每个 candidate 只输入一张原图和原问题，无 system prompt、无工具、无人工短答案指令；图像处理范围为 65,536--262,144 pixels。采样为 $T=1$、top-$p=1$、top-$k=-1$，每题 8 个独立 seed，总计 2,174,736 attempts。生成模型和语义评分模型是两个不同角色：
\begin{itemize}
  \item \textbf{生成/难度测量：}Qwen3-VL-8B-Instruct；
  \item \textbf{答案判定：}ArxivQA 使用 canonical A--Z rule；ThinkLite 使用 normalized exact/numeric/symbolic，必要时语义判定；VStar 使用 normalized exact 后语义判定；语义 judge 为本地 text-only Qwen2.5-72B-Instruct，temperature 0。
\end{itemize}

\subsection{保留规则与最终组成}

设 8 次中答对次数为 $k$。T1 规则为
\begin{equation}
\operatorname{decision}(k)=
\begin{cases}
\text{too hard}, & k=0,\\
\text{retain}, & 1\le k\le7,\\
\text{too easy}, & k=8.
\end{cases}
\end{equation}
截断、缺失或无法完成验证的 candidate 进入 unresolved；T2 明确不使用，T1 后不做按结果再平衡。最终 194 个 attempt 截断；candidate 决策为 too easy 168,378、too hard 25,762、retain 77,541、unresolved 161。

\begin{table}[t]
\centering
\caption{T1-retained-final-v2 的来源组成。}
\label{tab:data}
\small
\begin{tabular}{lrr}
\toprule
来源 & rows & share \\
\midrule
VStar & 39,205 & 50.56\% \\
ArxivQA & 25,393 & 32.75\% \\
ThinkLite & 12,943 & 16.69\% \\
\midrule
合计 & 77,541 & 100.00\% \\
\bottomrule
\end{tabular}
\end{table}

最终 samples/content SHA-256 的短标识分别为 \code{06e5b1b9...611a} 与 \code{5ab99622...2934}。formal 80-step schedule 每 step 固定 120 VStar、77 ArxivQA、59 ThinkLite，without replacement，seed 42；独立的 T1-PROBE256 使用 seed 20260807，组成 120/77/59，并从训练 schedule 中排除。当前 BS16 gate 只读取每个 formal batch 的小前缀，而非真的实现每 step 120/77/59；因此 step 间 source composition 仍会波动，单 step reward 不能画成严格 learning curve。

\paragraph{数据边界。}
CoreDev 只做了 exact image-byte leakage gate，尚无 perceptual near-duplicate gate。更重要的是，当前 DeepEyes-compatible source routing 将 ThinkLite 送入 single-turn/no-tool 文本 prompt；当前实现不把其图像放入该 prompt。对依赖图形的 ThinkLite 问题，这可能引入答案先验/文本 shortcut，是 formal scale 前必须单独量化的风险，而不是在总 reward 中被平均掉。

\section{Crop 协议：当前正在执行的路径}

\subsection{交互与像素语义}

VStar/ArxivQA 初始 context 为 system prompt 加 \{原图，question+user suffix\}。Policy 可输出 Hermes XML/JSON 工具调用；\code{bbox\_2d} 是 Qwen3-VL 原图相对 $0..1000$ 坐标。runtime 将其恰好一次转换为 immutable source-image pixel box，裁出 PIL RGB 图，再作为新的视觉 observation 送回 policy。每次 crop 都从原图裁，不对前一个 crop 递归裁；最多 6 次 observation。VStar 的 GT boxes 只用于 IoU/coverage audit，不进入 prompt、reward 或工具执行；ArxivQA 没有 GT box。

工具 observation 以新的 user/tool-response context 出现，内容为 crop image 加同一条 user suffix。它使下一次 assistant generation 看见新图，但本身不作为 policy action 训练。

\subsection{Crop 的 exact visible prompt}

当前 visible system/user 文本有意对齐 DeepEyes public V2。完整 literal 放在附录~\ref{app:prompts}；其语义可概括为：只暴露 \code{image\_zoom\_in\_tool}，要求模型先思考，需要时输出一个 bbox 工具调用，最后以 \code{<answer>...</answer>} 返回答案。公开 system schema 内保留了 \code{required:["bbox"]} 的原始 typo，但 runtime 可执行 schema 正确要求 \code{bbox\_2d}；这是 visible fidelity 与 executable safety 的显式分离。

ThinkLite 不暴露 Crop，单轮 suffix 为：
\begin{Verbatim}[breaklines=true,breakanywhere=true,fontsize=\scriptsize]
Let's think step by step and output the final answer within \boxed{}.
\end{Verbatim}

\section{当前 Crop reward 与 judge}

\subsection{Visual route：VStar/ArxivQA}

令 $A\in\{0,1\}$ 为答案是否正确；$F=0$ 表示格式合法，$F=-1$ 表示 think/answer 协议不合法；$C=1$ 当且仅当 $A=1$ 且至少一次 Crop runtime 执行成功。当前 reward 为
\begin{equation}
\boxed{\Rcrop^{\rm visual}=0.8A+0.2F+1.2C.}
\label{eq:crop-visual}
\end{equation}

\begin{table}[t]
\centering
\caption{Visual route 的典型 reward。}
\label{tab:crop-reward}
\small
\begin{tabularx}{\columnwidth}{Yc}
\toprule
轨迹结果 & Reward \\
\midrule
格式合法、直接答对 & 0.8 \\
格式合法、成功 Crop 后答对 & 2.0 \\
格式合法、答错 & 0.0 \\
格式非法、答错 & $-0.2$ \\
\bottomrule
\end{tabularx}
\end{table}

$A$ 并非由视觉模型重新看图判定。每一条 VStar/ArxivQA trajectory 都抽取最后一个 \code{<answer>}；没有 answer tag 时把整段回复交给 \textbf{text-only Qwen2.5-72B-Instruct}，输入只有 question、standard answer 和 model answer，判断语义是否一致。服务通过 OpenRouter 固定 DeepInfra provider；visual judge sampling 为 temperature 0.3、top-$p=1$、max 16 tokens。答案文本达到 1,000 characters 时强制 $A=0$。当前 resilient 配置把耗尽重试的瞬时网络失败局部化为该 trajectory 的零分；auth/model identity 错误仍应中止。

因此术语必须写准确：它是“visual-task final-answer text judge”，不是“看 Crop 是否正确的 VLM judge”。当前 reward 也没有直接奖励 bbox IoU、crop area、target quality 或因果 answer gain；这些量仅被记录为 audit metrics。

\subsection{ThinkLite route}

ThinkLite 不调用工具，reward 为
\begin{equation}
\boxed{\Rcrop^{\rm math}=1.2A+0.4F.}
\label{eq:crop-math}
\end{equation}
从最后一个 post-think \verb|\boxed{}| 抽取答案：先使用 deterministic \code{math\_verify}；规则未判为正确且存在候选答案时，再调用相同的 text-only Qwen2.5-72B-Instruct，temperature 0。没有 boxed answer 时 $A=0,F=-1$，且不发起语义 judge。

\subsection{Reward 的已知偏置}

式~\eqref{eq:crop-visual} 中“正确+Crop”的 2.0 是“正确+直接回答”0.8 的 2.5 倍。它能快速教会工具调用，但也把“曾成功执行 Crop”与“Crop 对答案真正有帮助”混合在一起。当前 gate 的首要问题是检验 DeepEyes fidelity 是否能在 full tuning/大 $n$ 下产生 held-out 正向信号；若工具率快速饱和而准确率不升，下一步应减少无条件 tool bonus 或改为 real-vs-sham counterfactual gain，而不是继续依据训练 reward 扩大规模。

\section{TGVF 协议与 reward 路线}

\subsection{Representation 与工具语义}

\TGVF{} arm 的 production representation 默认仍是 RP66 step-2000：Qwen3-VL-8B-Instruct contextual hidden state（layer $-1$）驱动冻结的 TGVF Adapter，输出 main $D$ 与 DeepStack 8/16/24；adapter SHA-256 为 \code{3429dc83880d...65ce5c9}。RP66 的 first-200 utility 已证明 image+$D^+$ 为 83.5\%，相对 image-only 72.5\% 的 raw gain 为 $+11.0\pp$，相对 zero/wrong-$D$ 的 controlled gain 为 $+7.5/+6.5\pp$。这证明“正确 $D$ 可作为原图的补充”，但不证明 free policy 会生成正确 target、会在正确时机调用、或端到端 RL 会提升。

\TGVF{} 工具为 \code{tgvf\_focus\_tool(target: string)}。target 必须是自包含的视觉查询，同时说明“看什么”和“要提取什么属性/关系”；禁止只写物体名，也禁止写入猜测的最终答案。成功 observation 的可见文本为
\begin{Verbatim}[breaklines=true,breakanywhere=true,fontsize=\scriptsize]
Focused visual observation for target:
"{target}"
\end{Verbatim}
并把 target-conditioned $D$ 作为新视觉信号加入后续 context。完整 TGVF system/user prompt 见附录~\ref{app:tgvf-prompt}。

RP67 的 image-axis loss 没有在 Stage1 accuracy 上形成明确 promotion，且曾观察到过强 donor-image rejection/norm shortcut 风险。因而 scale 研究的主 arm 应先固定 RP66；RP67 可作为只改变 representation 的一变量 follow-up，不能与 batch/full-tuning/reward 同时变化后再归因。

\subsection{Stage3-shaped reward 的准确公式}

未来完整 \TGVF{} reward 参考旧 Stage3 的有效语义，但使用当前 runtime 实现：
\begin{equation}
\boxed{\Rtgvf=2A_{\rm gated}+T+F_{\rm focus}+G_{\rm grounding}+P.}
\label{eq:tgvf}
\end{equation}

其五项定义如下。
\begin{enumerate}
  \item $A_{\rm gated}$：答案正确为 1；但 sidecar 标为 \code{needed} 且没有成功 \TGVF{} observation 时，答案项被 gate 为 0。它阻止“工具必要题上直接猜对仍拿满答案分”。
  \item $T$：工具决策 reward。sidecar 标签为 needed/optional/unnecessary；成功使用与不使用的 base score 分别为 $(+1,-2)$、$(+0.5,0)$、$(-0.5,+1)$，再乘 label confidence；第一次后的每次额外调用罚 $0.05$。
  \item $F_{\rm focus}$：target 是否与问题相关且聚焦，judge raw 0/1/2 映射为 $0/0.5/1$。
  \item $G_{\rm grounding}$：post-tool reasoning 是否被图像与 observation 支持，raw 0/1/2 映射为 $-1/0.5/1$。负分是防止 hallucinated evidence 的关键。
  \item $P$：任一 protocol error 时为 $-1$，否则为 0。
\end{enumerate}

\subsection{TGVF 到底用哪些 judge}

\begin{table*}[t]
\centering
\caption{训练 reward 中的 judge 边界。Evaluation judge 是另一套合同。}
\label{tab:judges}
\scriptsize
\begin{tabularx}{\textwidth}{lYYYl}
\toprule
路径 & 被判定对象 & 模型 & 模态/输入 & 是否在当前 Crop \\
\midrule
Crop visual $A$ & final answer 与 GT 是否一致 & Qwen2.5-72B-Instruct & text-only；Q/ref/candidate；每条 visual trajectory & 是 \\
Crop ThinkLite $A$ & math/开放答案等价 & rule-first + Qwen2.5-72B-Instruct & text-only fallback & 是 \\
TGVF answer $A$ & final answer correctness & rule-first + Qwen2.5-72B-Instruct & text-only semantic fallback & 否 \\
TGVF $F/G$ & target focus；post-tool grounding & Qwen3-VL-32B-Thinking & 原图+问题+target+reasoning+answer & 否；仅 full-shaped \\
CoreDev scoring & 各 benchmark 官方指标 & rules / 本地 Qwen2.5-72B-Instruct & VLMEvalKit task-specific & 后处理 \\
\bottomrule
\end{tabularx}
\end{table*}

PRL-09-R1 smoke 曾完整启用 32B visual judge，验证了 focus/grounding 调用链。随后本地 32B 与 4-GPU actor colocate 造成显存/headroom 与 HTTP 400 风险；PRL-09-R3 因此关闭 $F/G$，实际优化的是 $2A_{\rm gated}+T+P$，并未加载 32B。R3 的 80-step 累计工具尝试率为 99.93\%，mean answer accuracy 约 59.07\%，CoreDev 没有形成稳定全面增益。该结果说明 ATP/no-vision reward 容易学成“几乎必调工具”，不能作为完整式~\eqref{eq:tgvf} 已失败或已成功的证据。

若未来恢复 $F/G$，32B 可本地独占 GPU 或改用 API；但必须先做 latency/cost gate。当前 Crop 不应为此额外加载 32B，因为它的 reward 根本不需要视觉 judge。

\section{GRPO、policy loss 与 token ownership}

每个 prompt 的 $n=16$ 条 trajectory 构成一个 GRPO group。令 scalar reward 为 $r_i$，则
\begin{equation}
\hat A_i=\frac{r_i-\bar r}{s_r+10^{-6}}.
\end{equation}
同组 reward 全相同则 advantage 为 0。policy ratio 为
\begin{equation}
\rho_{it}=\exp\bigl(\log\pi_\theta(a_{it})-\log\pi_{\rm beh}(a_{it})\bigr),
\end{equation}
并使用 $\epsilon=0.2$ clipped PPO 与 dual clip 3。当前 loss mode 为 \code{deepeyes\_official\_micro\_token\_mean}：每个固定 local microbatch 先在有效 policy token 上取 mean，再在 gradient-accumulation micros/FSDP ranks 间平均。

\begin{figure}[t]
\centering
\fbox{\parbox{0.94\columnwidth}{\small
\textbf{Policy mask.}
进入 policy loss 的只有模型自己生成的 assistant tokens：reasoning、tool-call XML/JSON 与 final answer。初始 prompt、system/user template、原图 token、Crop/$D$ observation、tool-response wrapper 与 padding 的 mask 均为 0。工具输出会影响后续 action 的条件分布，但其 token 本身不被当成 policy action 反向传播。}}
\caption{工具 observation 与 policy loss 的关系。}
\label{fig:mask}
\end{figure}

这回答了“工具输出是否参与 policy loss”：\textbf{不直接参与}。不过 full tuning 下，后续 assistant action 的梯度会经过其所依赖的视觉计算图/重放输入更新 vision、projector 和 language parameters；这与把 observation token 自身纳入 action loss 是两件不同的事。

当前 KL 两条路径都关闭。日志中的 \code{rollout\_corr/kl}$\approx0.04$--$0.07$ 是 vLLM rollout logprob 与训练侧 logprob 的诊断差，不是 reference-policy KL penalty；\code{actor/ppo\_kl}=0 与该配置一致。

\section{当前 gate 的在线健康度}

\subsection{step 1--6 固定快照}

为避免运行结束时间改变报告口径，表~\ref{tab:online} 固定使用 2026-08-08 16:40 JST 已完整落盘的 step 1--6，共 1,536 trajectories。此后 step 7/8 与 T1-PROBE256 应在最终实验结论中另行追加。

\begin{table}[t]
\centering
\caption{当前 gate step 1--6 的 trajectory 汇总。}
\label{tab:online}
\small
\begin{tabularx}{\columnwidth}{Yc}
\toprule
指标 & 值 \\
\midrule
raw answer acc（混合训练 batch） & 51.50\% \\
mean scalar reward & 0.8285 \\
任意 Crop call / all trajectories & 68.03\% \\
任意成功 Crop / all trajectories & 67.06\% \\
任意 Crop call / visual trajectories & 87.08\% \\
mean Crop calls/trajectory & 1.726 \\
correct+successful-crop $C$ rate & 31.51\% \\
format error rate & 4.43\% \\
72B judge calls & 1,543 \\
judge failures / aborted trajectories & 0 / 0 \\
\bottomrule
\end{tabularx}
\end{table}

\begin{table}[t]
\centering
\caption{前 6 step 的 on-policy batch 指标；仅作健康性诊断。}
\label{tab:step-series}
\scriptsize
\begin{tabular}{lrrrrrr}
\toprule
step & 1 & 2 & 3 & 4 & 5 & 6 \\
\midrule
score & .648 & .914 & .952 & .787 & .800 & .870 \\
raw acc (\%) & 45.3 & 55.9 & 60.2 & 50.0 & 44.1 & 53.5 \\
\bottomrule
\end{tabular}
\end{table}

这些 batch 的 VStar/ArxivQA/ThinkLite 比例不同，且 rollout 随机，因此不能把高低顺序解释为单调学习曲线。step 5 pre-clip grad norm 曾到 42.62，但 grad clip 为 1，前后 steps 回到 3--12；entropy 有限，PPO clip fraction 为 0，未见 NaN、OOM、collapse 或全组 abort。

\subsection{速度与资源}

前 6 steps 分别用时 2,037、2,503、2,055、2,829、1,776、1,688 秒，平均 35.8 分钟/step。主要瓶颈是 old-logprob（704--1,268 秒）与 actor update（720--1,284 秒）；generation 为 90--208 秒，checkpoint 为 115--134 秒。72B reward judge 已并行批处理，除早期 transient latency 外不是总 wall-time 主瓶颈。训练 phase 峰值显存约 85.5 GiB allocated、93.5 GiB reserved/GPU；CPU snapshot 大部分 idle，当前不是 CPU bottleneck。

因此“GPU utilization 看起来阶段性稀疏”来自 colocated RL 的 rollout、logprob、update、save 和 weight-sync 相位切换，不等于整个 step 都在同一个 kernel 上低效。进一步加速优先级应是 old-logprob/update 的 sequence packing 与减少非必要 checkpoint，而不是增加 judge GPU。

\section{历史证据如何约束当前判断}

\begin{table*}[t]
\centering
\caption{旧 pilot、PRL-09 与当前 PRL-13 的可比边界。}
\label{tab:history}
\scriptsize
\begin{tabularx}{\textwidth}{lYYYY}
\toprule
实验 & Model/update & 数据与 rollout & 工具/reward & 可解释结论 \\
\midrule
旧 Thinking Pilot & Qwen3-VL-8B-Thinking；decoder LoRA r64；LR $10^{-5}$ & DeepEyes47K；BS16$n8$；80 & TGVF/Crop；0.8/0.2/1.2 & TGVF 未超过 baseline；Crop 坐标 bug，方法结论作废 \\
PRL-04/05/06 & Instruct；decoder LoRA；LR $10^{-6}$ & 早期 T1；BS16/32$n8$；20--80 & 修复 Crop；tool weight 1.2 或 0.2 & 更小 tool weight/BS32 未形成稳定 held-out 增益 \\
PRL-09-R3 & Instruct；decoder LoRA r64；RP66 frozen & mixed-T1-v2；BS16$n8$；80 & TGVF；$2A_g+T+P$，无 32B $F/G$ & 工具率饱和；不能代表完整 shaped reward \\
当前 PRL-13 gate & Instruct；\textbf{full vision+projector+LM}；LR $10^{-6}$ & corrected mixed-T1-v2；BS16$n16$；8；8 GPU & native Crop；DeepEyes official reward/judge & 首次隔离 full tuning+$n16$；待 step0/8 固定评测 \\
Formal scale（未启动） & 同 PRL-13 & BS256$n16$；20/80 & 先 Crop，后 matched TGVF & 只有小 gate 正向才扩展 \\
\bottomrule
\end{tabularx}
\end{table*}

旧实验没有稳定提升，不足以证明“mixed 数据必须达到官方全部规模才有效”；它只说明 BS16$n8$、decoder LoRA、强 tool bonus 的组合没有提供可靠信号。当前设计用一个可 scale 的小 gate 分开测试三个最可能的差异：$n$ 从 8 到 16、视觉路径从 frozen 到 full trainable、协议/reward 对齐 DeepEyes。batch 16 仍不足以验证 BS256 的统计稳定性，但比直接投入 80-step formal run 更适合先判断方向。

\section{评测计划与扩展门槛}

\subsection{当前 Crop gate}

step 8 保存完成后按以下顺序做评测：
\begin{enumerate}
  \item \textbf{T1-PROBE256：}内置、固定、训练排除；报告 overall/source accuracy、visual Crop rate、mean calls、format error、nonzero GRPO-group rate；
  \item \textbf{受控 step-0 vs step-8：}同一 native Crop prompt、sampling、pixel cap、工具上限与 scorer，只改变 checkpoint；
  \item \textbf{CoreDev-2511：}七项 VLMEvalKit；同时保留 direct/base policy 端到端 baseline，但明确 prompt/runtime 差异；
  \item \textbf{Grounding-200/VStar GT audit：}首框/最佳 IoU、GT coverage、crop area 与正确率的条件关系；它是工具健康指标，不并入 reward。
\end{enumerate}

promotion 不能只看 training reward 上升。建议小 gate 至少满足：无 crash/judge systematic failure；step-8 对 step-0 的固定 probe 不下降；CoreDev macro/关键 vision slice 有一致正向或至少不损害 reasoning；工具率不能以接近 100\% 的方式独立于题目必要性饱和。若 gate 不满足，先修数据路由/reward，不扩到 BS256。

\subsection{TGVF matched arm}

Crop 给出正向或可解释的 signal 后，\TGVF{} arm 应尽量只换工具：同 Qwen3-VL-Instruct、同 T1 schedule、同 prompt batch/$n$/LR/full-vs-LoRA 决策、同 answer judge、同 checkpoint/evaluation。第一主 arm 建议 RP66；RP67 作为独立 follow-up。Reward 分三层报告：
\begin{enumerate}
  \item ATP-only：$2A_{\rm gated}+T+P$，便宜但已知会 over-call；
  \item full-shaped：再加 32B $F/G$，验证 target/grounding；
  \item causal gate：真实 $D$ 对 zero/wrong/sham $D$ 的 answer gain，作为最终是否真正 foveate 的核心指标。
\end{enumerate}

如果本地 32B 影响 actor headroom，应把 VLM judge 放到独立 GPU/service 或 API，不与训练 GPU0 colocate。无论部署方式，visual judge failure 必须 sample-local 且显式计 coverage，不能用默认 0 冒充“judge 判为中性”。

\section{复现性与当前风险}

\begin{enumerate}
  \item \textbf{运行依赖已被清理：}当前进程的 cwd 指向已删除的 \code{tgvf-e2e-rl-prl13-integration}；train/probe schedule sentinel 也不在原位置。代码已加载入内存，所以当前 run 可继续，但若进程中断，\code{resume=auto} 不能在未重建这些文件前直接恢复。
  \item \textbf{身份来源：}W\&B metadata 记录 code commit \code{0a74c199...57672cfc}，但该 object 当前 main repo 不可 checkout。完整 live resolved config 仍保存在 Ray worker log，报告数值以它和 trajectory JSONL 为准，而不是以当前 HEAD 推测。
  \item \textbf{checkpoint 策略：}当前每步保存但只保留最新 actor；每个 live checkpoint 约 128 GiB。step 8 必须在任何清理前固定；若需要 step curve，应提前保留指定 step，而不是事后期待被删目录仍有权重。
  \item \textbf{W\&B 指标缺口：}当前 W\&B 有 scalar score/reward、长度、梯度和 timing，但没有 answer/tool/format component means。完整分解应从每步 256-row trajectory JSONL 重算。
  \item \textbf{ThinkLite route：}当前文本单轮路由可能丢失依赖图像的视觉信息；formal scale 前应做 image-present vs text-only probe，决定恢复图像还是剔除该 slice。
  \item \textbf{训练 loss fidelity：}当前 micro/GPU=32，而 DeepEyes 发布配置的 micro 更小；loss reduction 保持 micro-token-mean，但 micro partition 仍可能改变等权平均边界。报告应把它列为 capacity adaptation，而非数学完全等价。
\end{enumerate}

\section{结论与证据索引}

当前最准确的结论不是“RL 已有效”，也不是“RL 又失败”，而是：\textbf{PRL-13-A 已把此前最关键的 scale/vision-path 差异装进一个可承受的 8-step gate，当前机械与数值健康；能力结论尚欠固定 step-0/8 evaluation。}它使用正确的 Instruct T1-v2 77,541 条数据、native source-pixel Crop、full vision-policy update、$n=16$ 与 DeepEyes final-answer reward。当前 judge 是 text-only Qwen2.5-72B-Instruct，32B VLM 不在这条 Crop 训练路径中。

\TGVF{} 仍有明确研究依据：RP66 的 image+$D$ additive utility 已成立；真正未解决的是 free target/tool policy 与 grounding。未来 TGVF reward 应保留 needed-tool answer gate 和 negative grounding，同时避免把工具成功本身当作效用。完整 $F/G$ 才使用 Qwen3-VL-32B-Thinking；PRL-09-R3 的 no-vision 结果不能代替该实验。

\paragraph{关键证据路径。}
当前 trajectory/checkpoint root：
\begin{Verbatim}[breaklines=true,breakanywhere=true,fontsize=\scriptsize]
artifacts/policy/PRL-13-A-qwen3-instruct-grpo-bs256-n16-native-crop-t1-stratified-80step-gpu0123/pilot-bs16-n16-micro32-fused-flex-ws8/
\end{Verbatim}
其余证据路径如下；live resolved config 以 Ray session worker log 为准。
\begin{Verbatim}[breaklines=true,breakanywhere=true,fontsize=\scriptsize]
artifacts/data/policy_rl/
  T1-04-INSTRUCT-FULL-MIXED-T1-RETAINED-FINAL-v2/manifest.json
artifacts/data/policy_selection/t1/
  T1-04-QWEN3-INSTRUCT-512-FULLIMAGE-271842-GPU0123/
  run-config.canonical.json
src/tgvf_rl/protocol/tool_prompts.py
src/tgvf_rl/rewards/stage3_shaped.py
\end{Verbatim}

\clearpage
\onecolumn
\appendix
\small
\setlength{\parskip}{0pt}
\section{当前 Crop visible prompt（exact literal）}
\label{app:prompts}

\subsection{System prompt}

\begin{Verbatim}[breaklines=true,breakanywhere=true,fontsize=\scriptsize]
You are a helpful assistant.

# Tools
You may call one or more functions to assist with the user query.
You are provided with function signatures within <tools></tools> XML tags:
<tools>
{"type":"function","function":{"name":"image_zoom_in_tool","description":"Zoom in on a specific region of an image by cropping it based on a bounding box (bbox) and an optional object label.","parameters":{"type":"object","properties":{"bbox_2d":{"type":"array","items":{"type":"number"},"minItems":4,"maxItems":4,"description":"The bounding box of the region to zoom in, as [x1, y1, x2, y2], where (x1, y1) is the top-left corner and (x2, y2) is the bottom-right corner."},"label":{"type":"string","description":"The name or label of the object in the specified bounding box (optional)."}},"required":["bbox"]}}}
</tools>

# How to call a tool
Return a json object with function name and arguments within <tool_call></tool_call> XML tags:
<tool_call>
{"name": <function-name>, "arguments": <args-json-object>}
</tool_call>

**Example**:
<tool_call>
{"name": "image_zoom_in_tool", "arguments": {"bbox_2d": [10, 20, 100, 200], "label": "the apple on the desk"}}
</tool_call>
\end{Verbatim}

\subsection{User suffix and post-tool suffix}

\begin{Verbatim}[breaklines=true,breakanywhere=true,fontsize=\scriptsize]
Think first, call **image_zoom_in_tool** if needed, then answer. Format strictly as:  <think>...</think>  <tool_call>...</tool_call> (if tools needed)  <answer>...</answer>
\end{Verbatim}

初始 user content 顺序为 \code{[original image, question + suffix]}。成功工具 observation 包含 \code{[crop image, suffix]}，在当前 adapter 中按 DeepEyes 语义转换为 user \code{<tool\_response>...} continuation。visible schema 的 \code{required:["bbox"]} 仅为 public prompt fidelity；runtime registry 使用 \code{required:["bbox\_2d"]}。

\section{TGVF visible prompt（RP66/RP67 policy line）}
\label{app:tgvf-prompt}

\subsection{System prompt}

\begin{Verbatim}[breaklines=true,breakanywhere=true,fontsize=\scriptsize]
You are a visual reasoning assistant.

Use tgvf_focus_tool when additional target-conditioned visual evidence
is needed to answer the question.

The target must be a concise, self-contained visual query specifying
both what to inspect and what visual evidence or relation to obtain.
It may request an attribute, text reading, texture, state, count,
comparison, or spatial relation.

Do not provide only an object name, and do not include a guessed final
answer or answer-option value.

After receiving a tool result, continue reasoning. You may call the
tool again if more visual evidence is needed, up to four times. When
sufficient evidence is available, provide a concise final answer.

Example valid tool call:

<tool_call>
{"name":"tgvf_focus_tool","arguments":{"target":"the small circular gauge's needle position for reading its value"}}
</tool_call>
\end{Verbatim}

\subsection{Qwen3-Instruct shared user prompt}

\begin{Verbatim}[breaklines=true,breakanywhere=true,fontsize=\scriptsize]
<image>
{question}

Think first, call an available visual tool if needed, then answer.

Begin every assistant turn by generating one <think>...</think> block
yourself; the chat template does not add it. After </think>, either make
one native tool call if more visual evidence is needed or give only the
final answer.

After completing your reasoning, give only the final answer without explanation:
- For multiple-choice questions, give only the option letter.
- For mathematics questions, give only the final value or expression.
- For other questions, give only a concise answer.
\end{Verbatim}

\section{Judge 输入边界（便于审计）}

当前 Crop visual answer judge 的 system 为 \code{You are a helpful assistant.}，user prompt 给出 Question、Standard Answer 与 Model Answer，要求只输出 0/1 consistency judgement。它不包含原图、crop、bbox、checkpoint、arm 或 policy reward。ThinkLite math fallback 给出 question/reference/student final answer，要求只输出 TRUE/FALSE。两者 served model 均为 \code{qwen/qwen-2.5-72b-instruct}。

\TGVF{} full-shaped 的 visual-quality judge 与之不同：Qwen3-VL-32B-Thinking 接收原图，并分别输出 \code{focus\_score} 与 \code{grounding\_score}（0/1/2）。focus 只看 question、image 与 target；grounding 看 post-tool reasoning 是否被视觉事实支持。final answer correctness 仍由独立 72B text judge/rules 决定，避免一个 VLM verdict 同时控制所有 reward 分量。

\end{document}
